\pg{435} d) Os entes creados, embora possuam todos uma perfeição \emph{finita}, não á possuem todos no mesmo grau, mas um é \emph{mais} perfeito, e outro o é \emph{menos}. Isto mostra que a perfeição das creaturas é \emph{participada}; porque, se o não fosse, não haveria razão por que numa creatura a perfeição fosse \emph{maior}, e noutra \emph{menor}. Sendo \emph{participada}, a perfeição das creaturas deve derivar de um ente, e de um ente, que a possua pela \emph{sua propria essencia}. Ora a perfeição, que um ente possue pela sua \emph{propria essencia}, deve ser \emph{infinita}; porque, se fosse \emph{finita}, seria \emph{participada}. Logo existe um ente \emph{infinitamente} perfeito, de que derivaram, como de um oceano immenso, as muitas e diversas perfeições das creaturas. O ente \emph{infinitamente} perfeito é \textsc{Deus}. Logo existe \textsc{Deus} \footnote{Em poucas palavras: Os entes creados são todos \emph{finitos}. Ora a existencia dos entes \emph{finitos} suppõe necessariamente a existencia do ente \emph{infinito}. Logo existe o ente \emph{infinito}. O ente \emph{infinito} é \textsc{DEus}. Logo existe \textsc{Deus}.\par Objectam: A diversidade das perfeições, que se encontram nos seres finitos, não demonstra a existencia de um ente infinito.\par \emph{A} objecção já está refutada no texto. Entre os seres finitos, um é \emph{mais} ou \emph{menos} perfeito que o outro. Ora uma coisa diz-se \emph{mais} ou menos perfeita, emquanto \emph{mais} ou \emph{menos} se approxima de um ente \emph{absolutamente perfeito}: assim uma photographia diz-se \emph{mais} ou \emph{menos} perfeita, emquanto \emph{mais} ou menos se approxima do original, que é uma coisa \emph{absoluta}. Logo deve existir um ente \emph{absoluta} e \emph{infinitamente} perfeito, que seja a causa das varias e \emph{limitadas} perfeições, que se encontram nas creaturas. --- O \emph{relativo}, que se não refira ao \emph{absoluto}, é uma contradicção.}.

e) A intelligencia humana conhece verdades \emph{absolutas} e

\emph{necessarias}, tanto da ordem \emph{especulativa} como da \emph{pratica}. Taes verdades devem ter uma \emph{razão sufficiente}. Esta \emph{razão sufficiente} não se encontra nos \emph{seres creados}, nem na propria \emph{intelligencia humana}. Não se encontra nos \emph{seres creados}; pois estes, sendo \emph{relativos} e \emph{contingentes}, não podem produzir o \emph{absoluto} e O 